\section*{Chapitre \ref{ch:b3:hamiltonian-mechanics} --- Mécanique hamiltonienne}

\begin{solution}{exo:b3:hamiltonian-mechanics:1}
(a) $p = m\dot x$, $H = p\dot x - L = p^2/2m + \tfrac12 kx^2 = E_k +
E_p$. (b) $\dot x = p/m$, $\dot p = -kx$, d'où $\ddot x = -(k/m)x$.
(c) $H = E$ est l'ellipse $x^2/(2E/k) + p^2/(2mE) = 1$~: demi-axes
$\sqrt{2E/k}$ et $\sqrt{2mE}$. (d) Au point le plus à droite ($x > 0$,
$p = 0$), $\dot p = -kx < 0$~: le point descend --- dans le sens
horaire, toujours.
\end{solution}

\begin{solution}{exo:b3:hamiltonian-mechanics:2}
(a) $p = mR^2\dot\theta$~; $H = p^2/2mR^2 - \tfrac12 m\omega^2R^2
\sin^2\theta - mgR\cos\theta = p^2/2mR^2 + U_{\text{eff}}(\theta)$.
(b) Conservé ($H$ ne contient pas $t$ explicitement), mais c'est $h$,
non l'énergie mécanique~: le travail du moteur y manque. (c) Les
courbes de niveau sont $p = \pm\sqrt{2mR^2(H - U_{\text{eff}})}$~: pour
une rotation lente, des ovales autour de $(0, 0)$ et une \omterm{def:b3:hamiltonian-mechanics:portrait}{séparatrice} passant par le col en
$\theta = \pi$~; pour une rotation rapide, deux familles d'ovales
autour de $(\pm\theta_{\text{eq}}, 0)$. (d) Cas rapide~: centres en
$\pm\theta_{\text{eq}}$~; cols en $\theta = 0$ \emph{et} $\theta =
\pi$~; par $\theta = 0$ passe une \omterm{def:b3:hamiltonian-mechanics:portrait}{séparatrice} en forme de huit
entourant les deux centres (les \omterm{prop:b3:lagrangian-mechanics:modes}{petites oscillations} franchissent le
bas), et par $\pi$ la \omterm{def:b3:hamiltonian-mechanics:portrait}{séparatrice} extérieure, au-delà de laquelle la
perle circule par le haut.
\end{solution}

\begin{solution}{exo:b3:hamiltonian-mechanics:3}
(a) Un vol est l'arc $p = \pm\sqrt{2m^2g(z_{\max} - z)}$ --- une
parabole couchée, parcourue de $(0, +p_0)$ jusqu'à $(z_{\max}, 0)$
puis redescendue jusqu'à $(0, -p_0)$. (b) Le rebond envoie
$(0, -p_0)$ sur $(0, +p_0)$~: un segment vertical qui ferme la boucle. (c)
$\oint p\,\dd z = 2\int_0^{z_{\max}}m\sqrt{2g(z_{\max} - z)}\,\dd z =
\tfrac43 m\sqrt{2g}\,z_{\max}^{3/2}$, et $I$ en est le quotient par $2\pi$. (d)
$z_n = \big[3nh/(4m\sqrt{2g})\big]^{2/3}$~: pour un neutron, $z_1 =
(\num{6.7e-8})^{2/3} = \qty{17}{\micro m}$ --- le bon ordre~:
l'expérience de Grenoble a trouvé l'état quantique gravitationnel le
plus bas vers \qty{15}{\micro m} (le traitement exact, avec les vraies
fonctions d'onde, donne \qty{14}{\micro m}).
\end{solution}

\begin{solution}{exo:b3:hamiltonian-mechanics:4}
(a) $\{x^2, p\} = 2x$. (b) $\{xp, H\} = p^2/m - x\,E_p'(x) = 2E_k -
x\,E_p'$~: sur un mouvement lié la moyenne temporelle de $\dd(xp)/\dd t$
est nulle, donc $\langle 2E_k\rangle = \langle x\,E_p'\rangle$ --- le
théorème du viriel (pour $E_p \propto x^2$~: $\langle E_k\rangle = \langle
E_p\rangle$~; pour $\propto -1/r$~: $2\langle E_k\rangle = -\langle
E_p\rangle$). (c) $\{L_z, x\} = y$, $\{L_z, p_x\} = p_y$~: $L_z$
engendre les rotations des positions comme des moments. (d) Jacobi avec $h
= H$~: $\{\{f, g\}, H\} = -\{\{g, H\}, f\} - \{\{H, f\}, g\} = 0$.
\end{solution}

\begin{solution}{exo:b3:hamiltonian-mechanics:5}
(a) $E_{\text{sep}} = mg\ell$ (l'énergie de la position inversée au
repos)~; maximum $|p| = 2m\ell^2\omega_0$ en bas. (b) À partir de
$p^2/2m\ell^2 = mg\ell(1 + \cos\theta) = 2mg\ell\cos^2(\theta/2)$. (c)
$\dot\theta = 2\omega_0\cos(\theta/2)$ se sépare~: $\omega_0t =
\ln\tan(\theta/4 + \pi/4)$, qui diverge quand $\theta \to \pi$~: le
sommet est approché, jamais atteint. (d) Juste au-dessous de la \omterm{def:b3:hamiltonian-mechanics:portrait}{séparatrice}, le
mouvement la longe~: le pendule s'insinue au voisinage du col, s'y
attarde --- le logarithme --- et finit par retomber~; tout col ralentit
logarithmiquement les trajectoires, et c'est pourquoi un bâton mal
équilibré semble hésiter avant de tomber.
\end{solution}

\begin{solution}{exo:b3:hamiltonian-mechanics:6}
(a) $\dot{\vect r} = (\vect p - q\vect A)/m$ et $\dot p_x =
-\partial H/\partial x = (q/m)(\vect p - q\vect A)\cdot\partial_x\vect
A$, etc.~; en éliminant $\vect p$ on retrouve $m\ddot x = qB\dot y$,
$m\ddot y = -qB\dot x$. (b) $H = (\vect p - q\vect A)^2/2m = \tfrac12
m\vect v^{\,2}$~: l'énergie cinétique seule --- constante, la force
magnétique étant perpendiculaire à $\vect v$. (c) Avec $m\dot x = p_x +
\tfrac12 qBy$~: $\pi_x = p_x - \tfrac12 qBy = m\dot x - qBy$, dont la
conservation est la première équation du mouvement (de même pour $\pi_y$)~; dans
le crochet, deux termes seulement survivent~: $\{p_x, \tfrac12 qBx\} = -\tfrac12
qB$ et $\{-\tfrac12 qBy, p_y\} = -\tfrac12 qB$, au total $-qB$. ($\pi_x/qB$ et $\pi_y/qB$
sont, au signe près, les coordonnées du centre guide du cercle.) (d)
D'après \cref{exo:b3:hamiltonian-mechanics:4}(d), le crochet de deux
grandeurs conservées est conservé --- ici c'est la constante $-qB$,
conservée trivialement et qui n'apprend rien de neuf~: aucune
troisième grandeur n'apparaît.
\end{solution}

\begin{solution}{exo:b3:hamiltonian-mechanics:7}
(a) L'application $(q, p) \mapsto (q + pt/m, p)$ est un cisaillement~:
base et hauteur du rectangle sont inchangées, l'aire aussi
(déterminant $1$). (b) Dans les variables $(x\sqrt{m\omega}, p/\sqrt{m\omega})$ l'écoulement
est une rotation en bloc à la vitesse $\omega$~; les rotations
conservent l'aire, et le changement de variables a pour déterminant
$1$. (c) Le champ de vitesse
$(v, -\omega^2x - \gamma v)$ a pour divergence $-\gamma$~: toute aire
se contracte comme $\eu^{-\gamma t}$, en spiralant vers l'origine ---
impossible pour un écoulement \omterm{def:b3:hamiltonian-mechanics:hamiltonian}{hamiltonien}. (d) Elle passe dans les
\omterm{def:b3:lagrangian-mechanics:coordinates}{degrés de liberté} ignorés~: les molécules de l'air et les phonons du
fil, dont le volume dans l'\omterm{def:b3:hamiltonian-mechanics:hamiltonian}{espace des phases} croît au moins d'autant
--- le germe du second principe.
\end{solution}

\begin{solution}{exo:b3:hamiltonian-mechanics:8}
(a) $\dot r = p_r/m$, $\dot\varphi = p_\varphi/mr^2$, $\dot p_r =
p_\varphi^2/mr^3 - E_p'(r)$, $\dot p_\varphi = 0$. (b) $H$ ne contient pas
$\varphi$~: $\{p_\varphi, H\} = -\partial H/\partial\varphi =
0$~; c'est la conservation du moment cinétique. (c) $H_{\text{rad}} = p_r^2/2m +
U_{\text{eff}}(r)$ avec $U_{\text{eff}} = p_\varphi^2/2mr^2 + E_p(r)$.
(d) $U_{\text{eff}}' = 0$~: $r_{\text{c}} = p_\varphi^2/mk$, et $E =
U_{\text{eff}}(r_{\text{c}}) = -mk^2/2p_\varphi^2$.
\end{solution}

\begin{solution}{exo:b3:hamiltonian-mechanics:9}
(a) $\{L_x, L_y\} = \{yp_z - zp_y,\ zp_x - xp_z\}$~: les seuls
crochets canoniques non nuls donnent $yp_x\{p_z, z\} + xp_y\{z,
p_z\} = -yp_x + xp_y = L_z$. (b) La permutation circulaire $x \to y \to z \to
x$ donne les deux autres. (c) $\{L^2, L_z\} = 2L_x\{L_x, L_z\} +
2L_y\{L_y, L_z\} = -2L_xL_y + 2L_yL_x = 0$. (d) Pour $H = L^2/2J$,
chaque $\{L_i, H\} = 0$~: $\vect L$ est fixe. Avec $H = \sum L_i^2/2J_i$
et des $J_i$ inégaux~: $\{L_x, H\} = L_yL_z(1/J_y - 1/J_z) \neq 0$ ---
seuls $L^2$ et $H$ survivent, et le solide culbute (le théorème de la
raquette du chapitre sur le solide du volume de deuxième année vit ici).
\end{solution}

\begin{solution}{exo:b3:hamiltonian-mechanics:10}
(a) $\oint p\,\dd q = 2\pi I = 2\pi E/\omega = nh$ donne $E_n =
n\hbar\omega$. (b) Aller et retour à $|p|$ constant~: $2pa = nh$, $p =
nh/2a$, $E = p^2/2m = n^2h^2/8ma^2$ --- exactement le puits infini du
volume de deuxième année. (c) $\omega = \sqrt{k/\mu} = \qty{4.1e14}{rad/s}$,
$\hbar\omega = \qty{4.3e-20}{J} = \qty{0.27}{eV}$, $\lambda = hc/
\hbar\omega = \qty{4.6}{\micro m}$~: infrarouge moyen (la bande
fondamentale du CO, qui sert à tracer le gaz interstellaire). (d) Non~:
les différences de niveaux sont inchangées par la demi-quantité
commune. L'énergie de point zéro se voit dans les comparaisons qui
dépendent du niveau absolu --- décalages isotopiques des énergies de
dissociation (H$_2$ contre D$_2$), ou vibrations qui persistent au
zéro absolu dans les cristaux.
\end{solution}

\begin{solution}{exo:b3:hamiltonian-mechanics:11}
(a) $E = p^2/2m + \tfrac12 k(t)x^2$, donc le long d'un mouvement $\dot E =
\partial E/\partial t = \tfrac12\dot kx^2$. (b) Sur une période à $k$
pratiquement fixée, $\langle\tfrac12 kx^2\rangle = E/2$, donc
$\langle\dot E\rangle = \dot k\,E/2k$. (c) $\dd\ln E = \tfrac12\dd\ln
k$~; comme $\omega = \sqrt{k/m}$, on a aussi $\dd\ln\omega = \tfrac12\dd\ln k$~:
$E/\omega$ est invariant. (d) $\omega \propto \ell^{-1/2}$ croît
de $\sqrt2$, donc $E$ croît de $\sqrt2 \approx 1.41$. Avec $E = \tfrac12
mg\ell\theta_0^2$, $\theta_0 \propto \ell^{-3/4}$~: $\theta_0' =
5^\circ \times 2^{3/4} = 8.4^\circ$. L'énergie est fournie par qui
tire le fil~: la tension dépasse $mg\cos\theta$ en moyenne (terme
centrifuge), si bien que le hissage fournit un travail net positif.
\end{solution}

\begin{solution}{exo:b3:hamiltonian-mechanics:12}
(a) $\oint p\,\dd\theta = 2\int_{-\pi}^{\pi}2m\ell^2\omega_0
\cos(\theta/2)\,\dd\theta = 4m\ell^2\omega_0\big[2\sin(\theta/2)
\big]_{-\pi}^{\pi} = 16m\ell^2\omega_0$. (b) $\omega_0 =
\qty{9.9}{rad/s}$~: $16 \times \num{e-3} \times \num{e-2} \times 9.9 =
\qty{1.6e-3}{J.s}$, soit $N \approx \num{2.4e30}$ états~: le grain
quantique d'un pendule de laboratoire est trente ordres de grandeur
sous tout observable. (c) $16 \times \num{e-26} \times
\num{e-20} \times \num{e14} = \qty{1.6e-31}{J.s}$, soit $N \approx
240$~: une libration moléculaire ne loge que quelques centaines
d'états quantiques, et ses bas niveaux sont résolus un à un par la
spectroscopie. (d) La frontière est là où les aires de phase du
mouvement valent quelques $h$~: les molécules et au-dessous sont
quantiques~; tout ce qui est visible est classique.
\end{solution}

\begin{solution}{pb:b3:hamiltonian-mechanics:1}
\textbf{1.} $m_{\text{p}}/m_{\text{e}} = 1836$~: le proton bouge 1836
fois moins~; $H$ est celui de l'énoncé, dans le plan de l'orbite.
\textbf{2.} $\dot r = p_r/m_{\text{e}}$, $\dot\varphi =
p_\varphi/m_{\text{e}}r^2$, $\dot p_r = p_\varphi^2/m_{\text{e}}r^3 -
k/r^2$, $\dot p_\varphi = 0$.
\textbf{3.} $\varphi$ est absent de $H$~: $p_\varphi$, le moment
cinétique, est conservé.
\textbf{4.} $U_{\text{eff}} = p_\varphi^2/2m_{\text{e}}r^2 - k/r$~:
mur répulsif aux petits $r$, queue coulombienne aux grands $r$, un
minimum entre les deux.
\textbf{5.} $U_{\text{eff}}' = 0$ en $r_{\text{c}} =
p_\varphi^2/m_{\text{e}}k$.
\textbf{6.} $E = k/2r_{\text{c}} - k/r_{\text{c}} = -k/2r_{\text{c}} =
-m_{\text{e}}k^2/2p_\varphi^2$~; $E_p = -k/r_{\text{c}} = 2E$~: le
rapport du viriel de toute orbite circulaire en $1/r$.
\textbf{7.} $m_{\text{e}}v^2/r = k/r^2$~: $v = \sqrt{k/m_{\text{e}}r}$,
$f_{\text{orb}} = v/2\pi r = (1/2\pi)\sqrt{k/m_{\text{e}}r^3}$.
\textbf{8.} $k = \qty{2.30e-28}{J.m}$~: $v = \qty{2.2e6}{m/s}$ ($v/c =
0.0075$), $f_{\text{orb}} = \qty{7.2e15}{Hz}$, $E = -k/2r =
\qty{-2.3e-18}{J} = \qty{-14}{eV}$.
\textbf{9.} $p_\varphi = m_{\text{e}}vr = \qty{1.0e-34}{J.s} \approx
\hbar$~: une orbite atomique porte un moment cinétique de l'ordre de
$\hbar$ --- la constante de Planck est inscrite dans la taille des atomes.
\textbf{10.} Tout atome devrait s'effondrer en $\sim\qty{e-11}{s}$, sa
lumière balayant continûment vers les courtes longueurs d'onde. On
observe~: des atomes éternels, émettant des raies fines et fixes.
\textbf{11.} $1/\lambda = R_H(1/n^2 - 1/n'^2)$ écrit toute fréquence
observée comme une \emph{différence} de termes~: l'énergie s'échange
entre des niveaux fixes $-hcR_H/n^2$.
\textbf{12.} Une grandeur fixée à des valeurs universelles ne doit pas
dériver quand l'atome est doucement perturbé (champs, collisions,
changements lents)~; un \omterm{prop:b3:hamiltonian-mechanics:adiabatic}{invariant adiabatique} est précisément ce qui
reste fixe sous perturbation lente --- quantifions donc l'\omterm{def:b3:lagrangian-mechanics:action}{action}.
\textbf{13.} $p_\varphi$ est constant sur l'orbite~: $\oint
p_\varphi\,\dd\varphi = 2\pi p_\varphi = nh$, c'est-à-dire $p_\varphi =
n\hbar$.
\textbf{14.} $r_n = (n\hbar)^2/m_{\text{e}}k = n^2a_0$, $a_0 =
\hbar^2/m_{\text{e}}k = \qty{5.29e-11}{m} = \qty{52.9}{pm}$.
\textbf{15.} $E_n = -m_{\text{e}}k^2/2n^2\hbar^2 = -E_{\text{I}}/n^2$,
$E_{\text{I}} = m_{\text{e}}k^2/2\hbar^2 = \qty{2.18e-18}{J} =
\qty{13.6}{eV}$.
\textbf{16.} $v_1 = k/\hbar = \qty{2.19e6}{m/s}$~; $v_1/c = k/\hbar c =
1/137$~: la constante de structure fine $\alpha$, qui mesure à quel
point l'atome est non relativiste.
\textbf{17.} $h\nu = E_{n'} - E_n$ donne $1/\lambda =
(E_{\text{I}}/hc)(1/n^2 - 1/n'^2)$~: $R_H = E_{\text{I}}/hc =
\qty{1.10e7}{m^{-1}}$, en accord avec la valeur mesurée \qty{1.097e7}{m^{-1}} à
la précision de nos constantes.
\textbf{18.} $2 \to 1$~: \qty{122}{nm} (ultraviolet lointain, Lyman
$\alpha$)~; $3 \to 2$~: \qty{656}{nm}, la raie rouge visible qui colore
les nébuleuses en émission~; $\infty \to 2$~: \qty{365}{nm}, le bord
proche ultraviolet de la série de Balmer.
\textbf{19.} $k \to 2k$~: les rayons sont divisés par $2$, $E_{\text{I}} \to
4 \times \qty{13.6}{eV} = \qty{54.4}{eV}$~; la raie $2 \to 1$ tombe à
$\qty{122}{nm}/4 \approx \qty{30}{nm}$, dans l'ultraviolet extrême.
\textbf{20.} $f_{\text{orb}}(n) = 2E_{\text{I}}/hn^3$~;
$\nu_{n \to n-1} = (E_{\text{I}}/h)(2n - 1)/n^2(n-1)^2$. Rapport
$\nu/f_{\text{orb}} = n^3(2n-1)/2n^2(n-1)^2$~: $3$ à $n = 2$, $1.17$
à $n = 10$, $1.015$ à $n = 100$.
\textbf{21.} À la limite des grands nombres quantiques, les
prédictions quantiques doivent rejoindre les prédictions classiques
--- ici, la fréquence rayonnée rejoint la fréquence orbitale.
\textbf{22.} $n = 0$ signifie $p_\varphi = 0$~: une «~orbite
circulaire~» de rayon nul, l'électron sur le proton --- un tel
mouvement n'existe pas.
\textbf{23.} L'hydrogène mesuré a un état fondamental de moment
cinétique \emph{nul} (et non $\hbar$), et plusieurs états distincts
partageant le même $n$ (les formes $s$, $p$, $d$ de la chimie) ---
deux faits impossibles pour une unique orbite circulaire.
\textbf{24.} $a_0 \propto 1/m$~: $a_0^\mu = \qty{256}{fm}$~;
$E_{\text{I}}^\mu = 207 \times \qty{13.6}{eV} = \qty{2.8}{keV}$
(rayons X). Le muon orbite deux cents fois plus près --- en partie
\emph{à l'intérieur} de la charge nucléaire pour les éléments lourds
--- si bien que ses niveaux mesurent les rayons nucléaires (et ont
aiguisé l'énigme moderne de la taille du proton).
\textbf{25.} Un seul \omterm{prop:b3:hamiltonian-mechanics:adiabatic}{invariant adiabatique} égalé à $nh$ livre $a_0 =
\qty{52.9}{pm}$, $E_{\text{I}} = \qty{13.6}{eV}$, $R_H =
\qty{1.10e7}{m^{-1}}$ --- des nombres que la théorie quantique complète
du \cref{ch:b3:hydrogen-atom} gardera inchangés, tout en remplaçant
les orbites qui les ont produits.
\end{solution}
